% ===== PRÉAMBULE CANONIQUE — à coller VERBATIM en tête de CHAQUE .tex =====
\documentclass[11pt,a4paper]{article}
\usepackage[utf8]{inputenc}\usepackage[T1]{fontenc}\usepackage{lmodern}
\usepackage[french]{babel}
\usepackage{amsmath,amssymb,mathtools}\usepackage{icomma}
\usepackage{siunitx}\sisetup{locale=FR}  % \qty{}{}, \si{}, \ang{}
\usepackage{xcolor}\usepackage{colortbl}
\usepackage{tikz}\usepackage{pgfplots}\pgfplotsset{compat=1.18}
\usepackage[siunitx,europeanresistors]{circuitikz} % circuits (élec.)
\usepackage[most]{tcolorbox}\usepackage{enumitem}\usepackage{array,booktabs}
\usepackage[a4paper,margin=2.2cm]{geometry}
\usetikzlibrary{arrows.meta,calc,angles,quotes,positioning,patterns,%
  backgrounds,shapes.geometric,decorations.pathmorphing,decorations.markings,intersections}
\definecolor{primary}{RGB}{222,49,99}\definecolor{primarydark}{RGB}{150,22,55}
\definecolor{defcol}{RGB}{0,114,178}\definecolor{methcol}{RGB}{52,92,148}
\definecolor{excol}{RGB}{0,135,90}\definecolor{attncol}{RGB}{200,40,55}
\tcbset{boxbase/.style={enhanced,breakable,boxrule=0.9pt,arc=2.5pt,
  left=3mm,right=3mm,top=2.3mm,bottom=2.3mm,fonttitle=\bfseries,coltitle=white}}
% --- boîtes TUTORIEL (noms EXACTS, ne pas renommer) ---
\newtcolorbox{cle}[1][]{boxbase,colback=defcol!6,colframe=defcol,
  title={\textsc{À retenir}\ifx\relax#1\relax\relax\else\textmd{ : \itshape #1}\fi}}
\newtcolorbox{methode}[1][]{boxbase,colback=methcol!7,colframe=methcol,sharp corners,
  title={\textsc{Méthode}\ifx\relax#1\relax\relax\else\textmd{ --- \itshape #1}\fi}}
\newtcolorbox{exemplebox}[1][]{boxbase,colback=excol!5,colframe=excol,
  title={\textsc{Exemple}\ifx\relax#1\relax\relax\else\textmd{ : \itshape #1}\fi}}
\newtcolorbox{piege}[1][]{boxbase,colback=attncol!6,colframe=attncol,
  title={\textsc{Piège}\ifx\relax#1\relax\relax\else\textmd{ : \itshape #1}\fi}}
\newtcolorbox{remarque}[1][]{boxbase,colback=black!4,colframe=black!45,
  title={\textsc{Remarque}\ifx\relax#1\relax\relax\else\textmd{ : \itshape #1}\fi}}
% --- boîtes EXERCICE / CORRIGÉ (noms EXACTS, auto-comptées) ---
\newtcolorbox[auto counter]{exo}[1][]{boxbase,colback=primary!4,colframe=primary,
  title={\textsc{Exercice}~\thetcbcounter\ifx\relax#1\relax\relax\else\textmd{ --- \itshape #1}\fi}}
\newtcolorbox[auto counter]{corrige}[1][]{boxbase,colback=excol!4,colframe=excol,
  title={\textsc{Corrigé}~\thetcbcounter\ifx\relax#1\relax\relax\else\textmd{ --- \itshape #1}\fi}}
\newcommand{\vv}[1]{\vec{#1}}\newcommand{\ds}{\displaystyle}
\newcommand{\R}{\mathbb{R}}\newcommand{\N}{\mathbb{N}}\newcommand{\Z}{\mathbb{Z}}
% =========================================================================
\begin{document}
\sisetup{per-mode=symbol}

\begin{exo}[MCU : lesquelles sont fausses ?]
Un objet $A$ de masse $m_A$ tourne au bout d'un fil et fait un tour en \qty{0,25}{\second} (MCU).
Indique la ou les propositions \textbf{fausses} et corrige-les.
\begin{enumerate}[label=\Alph*.]
  \item L'accélération tangentielle de l'objet est nulle.
  \item Les vecteurs force centripète et vitesse linéaire ont toujours la même direction et le même sens.
  \item Le vecteur vitesse tangentielle est perpendiculaire au vecteur force centripète.
  \item La force centripète dépend de la masse, de la vitesse angulaire $\omega$ et du rayon $R$.
  \item Pour doubler la norme de la vitesse tangentielle, il faut doubler la force centripète.
  \item Si le fil casse, le mouvement suit une trajectoire rectiligne tangente au cercle.
\end{enumerate}
\end{exo}

\begin{exo}[rotation d'un disque]
Un disque plein tourne uniformément autour de son centre $O$. Trois points A, B, C sont alignés sur un
rayon (A proche du centre, C au bord). Réponds par \textbf{Vrai} ou \textbf{Faux} : une seule affirmation
est incorrecte.
\begin{center}
\begin{tikzpicture}[>=Stealth,scale=0.72]
  \fill[primary!12] (0,0) circle (3);
  \foreach \r in {1,2,3}{ \draw[black!40] (0,0) circle (\r); }
  \fill (0,0) circle (1.7pt) node[below right]{$O$};
  \foreach \r/\lab in {1/A,2/B,3/C}{
     \fill[primarydark] (90:\r) circle (2.3pt);
     \node[right,font=\small] at (90:\r){$\lab$};
     \draw[->,defcol,line width=0.9pt] (90:\r) -- ++(180:{0.32*\r+0.18});}
  \draw[->,primary] (3.5,-0.6) arc (-12:32:3.5);
  \node[primary,font=\small] at (4.0,0.8){sens};
\end{tikzpicture}
\end{center}
\begin{enumerate}[label=\Alph*.]
  \item La vitesse tangentielle d'un point est nulle s'il est placé exactement au centre $O$.
  \item La vitesse angulaire d'un point est nulle lorsqu'il est placé au centre $O$.
  \item Toutes les parties du disque tournent à la même vitesse angulaire.
  \item Un point du bord (C) a une vitesse linéaire plus grande qu'un point proche du centre (A).
\end{enumerate}
\end{exo}

\begin{exo}[la fronde : capstone]
On fait tourner une pierre de masse $m=\qty{0,3}{\kilogram}$ au bout d'un fil de longueur
$R=\qty{0,4}{\metre}$, à raison de \num{2} tours par seconde ($\pi\approx\num{3,14}$).
\begin{enumerate}[label=\alph*)]
  \item Calcule la période $T$ et la vitesse angulaire $\omega$.
  \item Calcule la vitesse linéaire $v$.
  \item Calcule l'accélération centripète $a_c$, puis la force centripète $F_c$ (tension du fil).
  \item Le fil casse. Décris la trajectoire ultérieure de la pierre et nomme la loi invoquée.
\end{enumerate}
\end{exo}

\begin{exo}[le virage d'une voiture]
Une voiture de masse $m$ prend un virage circulaire de rayon $R$ à la vitesse $v$. C'est le frottement
des pneus qui fournit la force centripète $F_c=\dfrac{mv^2}{R}$.
\begin{enumerate}[label=\alph*)]
  \item Si la voiture double sa vitesse (même $R$), par quel facteur la force centripète nécessaire est-elle multipliée ?
  \item Si le rayon du virage est divisé par deux (même vitesse), par quel facteur $F_c$ est-elle multipliée ?
\end{enumerate}
\end{exo}
\end{document}
